\ifx\allfiles\undefined
\documentclass[12pt, a4paper, oneside, UTF8]{ctexbook}
\def\path{../config}
\input{\path/package.tex}

% 定理环境设置
\input{\path/theorem1_zh.tex}

\input{\path/custom.tex}

% 修改参数改变封面样式，0 默认原始封面、内置其他1、2、3种封面样式
\def\myIndex{3}


\ifnum\myIndex>0
    \input{\path/cover_package_\myIndex} 
\fi

\def\myTitle{高等数学笔记}
\def\myAuthor{M.K.}
\def\myDateCover{\today}
\def\myDateForeword{\today}
\def\myForeword{前言}
\def\myForewordText{
    
    这是一个基于\LaTeX{}模板的高数笔记．
    内容参考了包括但不限于以下资料：
    \begin{itemize}
        \item 《高等数学》（同济大学出版社）第七版
        \item 《吉米多维奇高等数学习题精选精讲》（山东科学技术出版社）第二版
        \item 《高等数学作业集》（哈尔滨工业大学出版社）第一版
        \item \href{https://space.bilibili.com/1035929235}{B站UP主：一高数}的高数课程
        \item \href{https://space.bilibili.com/42428180}{B站UP主：考研竞赛数学}的试题讲解与高数课程
        \item \href{https://space.bilibili.com/391930545}{B站UP主：Maki的完美算数教室}的高数内容

    \end{itemize}

    封面来源于\href{https://www.pixiv.net/users/3952}{おにねこ}的作品\href{https://www.pixiv.net/artworks/136551599}{雨雫とプレアデス}
    如有侵权请联系我删除．
}
\def\mySubheading{副标题}


\begin{document}
% \input{../config/cover}
\else
\fi
\chapter{无穷级数}
\section{正项级数及其收敛判定}
\begin{theorem}[比较判别法的极限形式]\label{thm:lim-comp}
    设 $\sum u_n$ 和 $\sum v_n$ 为正项级数，且 $\lim\limits_{n \to \infty} \frac{u_n}{v_n} = l$：
    \begin{itemize}
        \item 若 $0 < l < +\infty$，则 $\sum u_n$ 与 $\sum v_n$ 同敛散；
        \item 若 $l = 0$，且 $\sum v_n$ 收敛，则 $\sum u_n$ 收敛；
        \item 若 $l = +\infty$，且 $\sum v_n$ 发散，则 $\sum u_n$ 发散．
    \end{itemize}
\end{theorem}

\begin{theorem}[比值判别法 (D'Alembert)]
    若 $\lim\limits_{n \to \infty} \dfrac{u_{n+1}}{u_n} = \rho$：
    \begin{itemize}
        \item $\rho < 1$ 时级数收敛；
        \item $\rho > 1$ 时级数发散；
        \item $\rho = 1$ 时失效．
    \end{itemize}
\end{theorem}

\begin{theorem}[根值判别法 (Cauchy)]
    若 $\lim\limits_{n \to \infty} \sqrt[n]{u_n} = \rho$：
    \begin{itemize}
        \item $\rho < 1$ 时级数收敛；
        \item $\rho > 1$ 时级数发散；
        \item $\rho = 1$ 时失效．
    \end{itemize}
\end{theorem}

\begin{theorem}[积分判别法]
    设 $f(x)$ 在 $[M,+\infty)$ 上连续、恒正且单调递减，其中 $M \in \mathbb{Z}^+$，令 $a_n = f(n)$．则级数 $\displaystyle\sum_{n=M}^\infty a_n$ 与反常积分 $\displaystyle\int_M^{+\infty} f(x)\d x$ 同敛散．
\end{theorem}

\begin{corollary}[$p$-级数的敛散性]
    级数 $\displaystyle\sum_{n=1}^\infty \frac{1}{n^p}$ 当 $p > 1$ 时收敛，当 $p \le 1$ 时发散．
\end{corollary}
\begin{proof}
    当 $p \le 0$ 时，通项 $\dfrac{1}{n^p} \not\to 0$，级数发散．
    当 $p > 0$ 时，函数 $f(x) = \dfrac{1}{x^p}$ 在 $[1,+\infty)$ 上连续、恒正且单调递减，由积分判别法，级数 $\sum \dfrac{1}{n^p}$ 与反常积分 $\displaystyle\int_1^{+\infty} \dfrac{1}{x^p}\d x$ 同敛散．计算该积分：
    \[ \int_1^{+\infty} \frac{1}{x^p}\d x = 
        \begin{cases}
            \displaystyle\lim_{b\to+\infty}\frac{x^{1-p}}{1-p}\Big|_1^b = +\infty, & p < 1 \\[8pt]
            \displaystyle\lim_{b\to+\infty}\ln x\,\big|_1^b = +\infty, & p = 1 \\[8pt]
            \displaystyle\frac{1}{p-1}, & p > 1
        \end{cases} \]
    故当 $0 < p \le 1$ 时积分发散，$p > 1$ 时积分收敛．综上所述，$\sum \dfrac{1}{n^p}$ 收敛当且仅当 $p > 1$．
\end{proof}

\begin{corollary}[广义 $p$-级数]
    级数 \[\sum_{n=2}^\infty \frac{1}{n^\alpha (\ln n)^\beta}\] 的收敛性如下：
    \begin{itemize}
        \item 当 $\alpha > 1$ 时，对任意 $\beta$，级数收敛；
        \item 当 $\alpha < 1$ 时，对任意 $\beta$，级数发散；
        \item 当 $\alpha = 1$ 时
            \begin{itemize}
                \item $\beta > 1$ 时级数收敛；
                \item $\beta \le 1$ 时级数发散．
            \end{itemize}
    \end{itemize}
\end{corollary}

\begin{proof}[广义 $p$-级数的证明]
    \begin{enumerate}
        \item 当 $\alpha > 1$ 时，取实数 $p$ 使得 $1 < p < \alpha$．由极限关系
        \[ \lim_{n \to \infty} \frac{ \frac{1}{n^\alpha (\ln n)^\beta} }{ \frac{1}{n^p} } = \lim_{n \to \infty} \frac{1}{n^{\alpha-p} (\ln n)^\beta} = 0 \]
        可知存在充分大的 $N$，当 $n > N$ 时有 $\dfrac{1}{n^\alpha (\ln n)^\beta} \le \dfrac{1}{n^p}$．因 $p>1$，由收敛的最基本 $p$-级数结论和比较判别法知，原级数收敛．
        \item 当 $\alpha < 1$ 时，取实数 $p$ 使得 $\alpha < p < 1$．由极限关系（利用洛必达法则）
        \[ \lim_{n \to \infty} \frac{ \frac{1}{n^\alpha (\ln n)^\beta} }{ \frac{1}{n^p} } = \lim_{n \to \infty} \frac{n^{p-\alpha}}{(\ln n)^\beta} = +\infty \]
        可知对于充分大的 $n$，有 $\dfrac{1}{n^\alpha (\ln n)^\beta} \ge \dfrac{1}{n^p}$．因 $p<1$，级数 $\sum \frac{1}{n^p}$ 发散，故由比较判别法知，原级数发散．
        \item 当 $\alpha = 1$ 时，利用积分判别法．考察函数 $f(x) = \dfrac{1}{x (\ln x)^\beta}$，当 $x$ 充分大时 $f(x)$ 连续、恒正且单调递减．计算反常积分：
        \[ \int_2^{+\infty} \frac{1}{x (\ln x)^\beta} \d x = \int_{\ln 2}^{+\infty} \frac{1}{u^\beta} \d u \quad (\text{令 } u = \ln x) \]
        由反常积分 $p$-积分的判定结论知，该积分当且仅当 $\beta > 1$ 时收敛，$\beta \le 1$ 时发散．因此原级数与积分敛散性一致．
    \end{enumerate}
\end{proof}

\begin{example}
    若$\displaystyle\sum_{n=1}^{\infty}u_n$收敛，且$\lim\limits_{n\to\infty}\dfrac{v_n}{u_n}=1$，判定$\displaystyle\sum_{n=1}^{\infty}v_n$的敛散性．
\end{example}
\begin{solution}
    \textbf{无法判定} $\sum v_n$ 的敛散性．\\
    根据\hyperref[thm:lim-comp]{比较判别法的极限形式}，只有$\sum u_n$ 和 $\sum v_n$ 为正项级数时，可判断 $\sum v_n$ 收敛．

    \textbf{反例}：取
    \[ u_n = \frac{(-1)^n}{\sqrt{n}},\qquad v_n = \frac{(-1)^n}{\sqrt{n}} + \frac{1}{n} \]
    则 $\sum u_n$ 收敛，且
    \[ \frac{v_n}{u_n} = 1 + \frac{(-1)^n}{\sqrt{n}} \to 1 \]
    但 $\displaystyle\sum_{n=1}^{\infty} \frac{1}{n}$ 发散，故 $\sum v_n$ 作为收敛级数与发散级数之和，\textbf{发散}．
\end{solution}

\begin{example}
    判定级数$\displaystyle\sum_{n=1}^{\infty}\frac{2n-1}{2^n}$的敛散性．
\end{example}
\begin{solution}
    \textbf{法一（比值判别法）：} 设 $u_n = \frac{2n-1}{2^n}$，则
    \[ \lim_{n \to \infty} \frac{u_{n+1}}{u_n} = \lim_{n \to \infty} \frac{2(n+1)-1}{2^{n+1}} \cdot \frac{2^n}{2n-1} = \lim_{n \to \infty} \frac{2n+1}{2(2n-1)} = \frac{1}{2} < 1 \]
    由比值判别法知，级数收敛．

    \textbf{法二（比较判别法的极限形式）：} 选取收敛的几何级数 $\sum \left(\frac{3}{4}\right)^n$．
    \[ \lim_{n \to \infty} \frac{u_n}{(3/4)^n} = \lim_{n \to \infty} \frac{2n-1}{2^n} \cdot \frac{4^n}{3^n} = \lim_{n \to \infty} (2n-1) \left(\frac{2}{3}\right)^n = 0 \]
    由比较判别法的极限形式知，级数收敛．

    \textbf{法三（根值判别法）：}
    \[ \lim_{n \to \infty} \sqrt[n]{u_n} = \lim_{n \to \infty} \frac{\sqrt[n]{2n-1}}{2} = \frac{1}{2} < 1 \]
    （利用 $\lim_{n \to \infty} \sqrt[n]{n} = 1$ 可推得 $\lim_{n \to \infty} \sqrt[n]{2n-1} = 1$）．由根值判别法知，级数收敛．
\end{solution}

\section{交错级数与绝对收敛}
\begin{theorem}[Leibniz 判别法]
    若交错级数 $\sum_{n=1}^\infty (-1)^{n-1} u_n$ ($u_n > 0$) 在$n$充分大时满足：
    \begin{enumerate}
        \item $u_n \ge u_{n+1}$ (单调不增)；
        \item $\lim_{n \to \infty} u_n = 0$；
    \end{enumerate}
    则该级数收敛，且其余项 $|R_n| \le u_{n+1}$．
\end{theorem}

\begin{definition}[绝对收敛与条件收敛]
    若 $\sum |u_n|$ 收敛，称 $\sum u_n$ \textbf{绝对收敛}；若 $\sum u_n$ 收敛但 $\sum |u_n|$ 发散，称 $\sum u_n$ \textbf{条件收敛}．
\end{definition}

\begin{example}
    设$u_n=(-1)^n\ln(1+\dfrac{1}{\sqrt{n}})$，判定$\displaystyle\sum_{n=1}^{\infty}u_n$的敛散性．
\end{example}
\begin{solution}
    令 $v_n = |u_n| = \ln(1+\dfrac{1}{\sqrt{n}})$．
    \begin{enumerate}
        \item 显然随着 $n \to \infty$，$\frac{1}{\sqrt{n}} \to 0$，则 $v_n \to 0$．且函数 $y = \ln(1+x^{-\frac{1}{2}})$ 随 $x$ 增大严格单调递减，说明 $v_n$ 单调递减趋向于 $0$．由 Leibniz 判别法知，交错级数 $\sum u_n$ 收敛．
        \item 对于其绝对值级数 $\sum |u_n|$，由等价无穷小可知当 $n \to \infty$ 时，$v_n \sim \dfrac{1}{\sqrt{n}}$．由于 $p$-级数 $\sum \dfrac{1}{\sqrt{n}}$ 发散（$p = \frac{1}{2} \le 1$），由极限形式比较判别法知 $\sum |u_n|$ 发散．
    \end{enumerate}
    综上所述，该级数条件收敛．
\end{solution}

\section{幂级数}
\subsection{收敛半径与收敛域}
\begin{theorem}[Abel 定理]
    若 $\sum a_n x^n$ 在 $x_0 \neq 0$ 处收敛，则对于满足 $|x| < |x_0|$ 的一切 $x$ 都绝对收敛；若在 $x_0$ 处发散，则对于满足 $|x| > |x_0|$ 的一切 $x$ 都发散．
\end{theorem}
\begin{itemize}
    \item \textbf{收敛半径} $R$ 可通过以下公式求得：
        \begin{itemize}
            \item \textbf{比值法}：$R = \lim\limits_{n \to \infty} \left| \dfrac{a_n}{a_{n+1}} \right|$
            \item \textbf{根值法}：$R = \lim\limits_{n \to \infty}\dfrac{1}{\sqrt[n]{|a_n|}}$
        \end{itemize}
\end{itemize}

\begin{example}
    已知幂级数 $\displaystyle\sum_{n=1}^{\infty}a_n(x-1)^n$ 在 $x=2$ 处条件收敛，则改幂级数（~~~~~~）
    \begin{tasks}[style = itemize, label=\Alph*.](2)
        \task 收敛半径为 $2$
        \task 收敛区间为 $(0, 2]$
        \task 收敛域为 $(0, 2]$
        \task 收敛区间为 $(0, 2)$
    \end{tasks}
\end{example}
\begin{solution}
    幂级数 $\displaystyle\sum_{n=1}^{\infty}a_n(x-1)^n$ 的中心为 $x=1$．在 $x=2$ 处条件收敛，即 $|2-1|=1$ 处条件收敛．

    由 Abel 定理，在收敛区间内部幂级数绝对收敛，在边界上可能条件收敛或发散．
    若收敛半径 $R>1$，则 $x=2$ 在收敛区间内部，应绝对收敛，与条件收敛矛盾．故 $R \le 1$．
    又 $x=2$ 处收敛，故 $R \ge 1$．因此 $R=1$．

    收敛区间（开区间）为 $|x-1|<1$，即 $(0,2)$．

    在右端点 $x=2$ 处条件收敛，左端点 $x=0$ 处敛散性未知，例如\begin{enumerate}[label=\roman*.]
        \item 取$a_n = \begin{cases}
            \dfrac{(-1)^{\lfloor{\frac{n}{2}}\rfloor}}{n}, & n \text{ 为奇数} \\
            0, & n \text{ 为偶数}
            \end{cases}$
        ，则 $\sum a_n(x-1)^n$ 在 $x=0$ 处收敛；
        \item 取$a_n = \dfrac{(-1)^n}{n}$，则 $\sum a_n(x-1)^n$ 在 $x=0$ 处发散．
    \end{enumerate}
    
    

    因此 $\boxed{\text{D}}$ 正确．
\end{solution}

\begin{example}
     求下列幂级数的收敛域
     \begin{tasks}[label=(\arabic*)](3)
        \task $\displaystyle\sum_{n=1}^{\infty}\left(\frac{3^n}{n^2}\right)x^n$
        \task $\displaystyle\sum_{n=0}^{\infty}\frac{n^2+1}{n!}x^n$
        \task $\displaystyle\sum_{n=1}^{\infty}\frac{4^nx^{2n-1}}{n}$
    \end{tasks}
\end{example}
\begin{solution}
    \begin{enumerate}[label=(\arabic*)]
        \item 记 $a_n = \dfrac{3^n}{n^2}$，则收敛半径
        \[ R = \lim_{n\to\infty}\left|\frac{a_n}{a_{n+1}}\right| = \lim_{n\to\infty} \frac{3^n}{n^2}\cdot\frac{(n+1)^2}{3^{n+1}} = \frac13 \lim_{n\to\infty}\frac{(n+1)^2}{n^2} = \frac13 \]
        端点处：
        \begin{itemize}
            \item $x = \dfrac13$：$\displaystyle\sum_{n=1}^\infty \frac{3^n}{n^2}\cdot\frac{1}{3^n} = \sum_{n=1}^\infty \frac{1}{n^2}$ 收敛（$p$-级数，$p=2>1$）．
            \item $x = -\dfrac13$：$\displaystyle\sum_{n=1}^\infty \frac{(-1)^n}{n^2}$ 绝对收敛．
        \end{itemize}
        故收敛域为 $\boxed{\left[-\dfrac13,\;\dfrac13\right]}$．

        \item 记 $a_n = \dfrac{n^2+1}{n!}$，则收敛半径
        \[ R = \lim_{n\to\infty}\left|\frac{a_n}{a_{n+1}}\right| = \lim_{n\to\infty} \frac{n^2+1}{n!}\cdot\frac{(n+1)!}{(n+1)^2+1} = \lim_{n\to\infty} \frac{(n^2+1)(n+1)}{n^2+2n+2} = +\infty \]
        故收敛域为 $\boxed{(-\infty,\;+\infty)}$．

        \item 注意到该级数不是标准形式 $\sum a_n x^n$
        \[\rho = \lim_{n\to\infty}\left|\frac{a_{n+1}}{a_n}\right| = 4 \implies R = \sqrt{\frac{1}{\rho}} = \frac12\]
        端点处：$x=\pm\dfrac12$ 时，$\displaystyle\sum_{n=1}^\infty \frac{4^n}{n}\left(\pm\dfrac12\right)^{2n-1} = \pm 2\sum_{n=1}^\infty \frac{1}{n}$ 发散．
        故收敛域为 $\boxed{\left(-\dfrac12,\;\dfrac12\right)}$．
    \end{enumerate}
\end{solution}

\begin{theorem}[幂级数的分析性质 (逐项求导与积分)]
    设幂级数 $\displaystyle\sum_{n=0}^\infty a_n x^n$ 的收敛半径为 $R > 0$，其和函数 $S(x) = \displaystyle\sum_{n=0}^\infty a_n x^n$：
    \begin{itemize}
        \item \textbf{逐项求导定理}：$S(x)$ 在 $(-R, R)$ 内无限次可导，且逐项求导后：
        \[ S'(x) = \sum_{n=1}^\infty n a_n x^{n-1} \]
        逐项求导所得级数的收敛半径仍为 $R$．
        \item \textbf{逐项积分定理}：$S(x)$ 在 $(-R, R)$ 内连续可积，且：
        \[ \int_0^x S(t) \d t = \sum_{n=0}^\infty \frac{a_n}{n+1} x^{n+1} \]
        逐项积分所得级数的收敛半径仍为 $R$．
    \end{itemize}
    注意：尽管收敛半径不变，但在端点处的敛散性可能不同．
\end{theorem}

\subsection{和函数与幂级数的转换}
\begin{table}[H]
    \centering
    \caption{常用函数的幂级数展开式}
    \renewcommand\arraystretch{1.5}
    \begin{tabular}{lll}
        \hline
        \textbf{函数表达式} & \textbf{幂级数展开} & \textbf{收敛域} \\
        \hline
        $e^x$ & $\displaystyle\sum_{n=0}^\infty \frac{x^n}{n!}$ & $ x \in (-\infty, +\infty)$ \\
        $\sin x$ & $\displaystyle\sum_{n=0}^\infty (-1)^n \frac{x^{2n+1}}{(2n+1)!}$ & $ x \in (-\infty, +\infty)$ \\
        $\cos x$ & $\displaystyle\sum_{n=0}^\infty (-1)^n \frac{x^{2n}}{(2n)!}$ & $ x \in (-\infty, +\infty)$ \\
        $\dfrac{1}{1-x}$ & $\displaystyle\sum_{n=0}^\infty x^n$ & $ x \in (-1, 1)$ \\
        $\dfrac{1}{1+x}$ & $\displaystyle\sum_{n=0}^\infty (-1)^n x^n$ & $ x \in (-1, 1)$ \\
        $\ln(1+x)$ & $\displaystyle\sum_{n=1}^\infty (-1)^{n-1} \frac{x^n}{n}$ & $x \in (-1, 1]$ \\
        $-\ln(1-x)$ & $\displaystyle\sum_{n=1}^\infty \frac{x^n}{n}$ & $x \in [-1, 1)$ \\
        $\arctan x$ & $\displaystyle\sum_{n=0}^\infty (-1)^n \frac{x^{2n+1}}{2n+1}$ & $ x \in [-1, 1]$ \\
        $\dfrac{1}{2}\ln\left(\dfrac{1+x}{1-x}\right)$ & $\displaystyle\sum_{n=0}^{\infty}\frac{x^{2n+1}}{2n+1}$ & $x\in (-1,1)$\\
        \hline
    \end{tabular}
\end{table}

\begin{example}
    计算$\displaystyle\sum_{n=0}^{\infty}(2n+1)x^n$的收敛域及和函数，并求$\displaystyle\sum_{n=0}^{\infty}(-1)^n\frac{2n+1}{2^n}$
\end{example}
\begin{solution}
    令 \[S(x) = \sum_{n=0}^\infty (2n+1)x^n = 2\sum_{n=0}^\infty nx^n + \sum_{n=0}^\infty x^n\]
    已知 $\displaystyle\sum_{n=0}^\infty x^n = \frac{1}{1-x}$，收敛域为 $(-1, 1)$．\\
    利用逐项求导：\[\sum_{n=0}^\infty nx^n = x \sum_{n=1}^\infty nx^{n-1} = x \left(\sum_{n=0}^\infty x^n\right)' = x \left(\frac{1}{1-x}\right)' = \frac{x}{(1-x)^2}\]
    因此，和函数 $S(x) = \dfrac{2x}{(1-x)^2} + \dfrac{1}{1-x} = \boxed{\dfrac{1+x}{(1-x)^2}}$，收敛域为 $\boxed{(-1, 1)}$．\\
    代入$x=-\dfrac{1}{2}$，得$\displaystyle\sum_{n=0}^{\infty}(-1)^n\frac{2n+1}{2^n}=\boxed{\frac{2}{9}}$
\end{solution}

\begin{example}
    求下列幂级数的和函数$S(x)$
    \begin{tasks}[label=(\arabic*)](3)
        \task $\displaystyle\sum_{n=1}^{\infty}\frac{x^{n+1}}{n(n+1)}$
        \task $\displaystyle\sum_{n=1}^{\infty}\frac{1}{2n}x^{2n+1}$
        \task $\displaystyle\sum_{n=0}^{\infty}\frac{x^{3n}}{(3n)!}$
    \end{tasks}
\end{example}
\begin{solution}
    \begin{enumerate}[label=(\arabic*)]
        \item 易知级数收敛半径 $R=1$．当 $x=1$ 时，级数为 $\sum \frac{1}{n(n+1)}$ 收敛；当 $x=-1$ 时，级数为 $\sum \frac{(-1)^{n+1}}{n(n+1)}$ 绝对收敛．故收敛域为 $[-1, 1]$．\\
        当 $x \in (-1, 1)$ 时，和函数 $S(x)$ 满足：
        \[S''(x) = \sum_{n=1}^\infty x^{n-1} = \frac{1}{1-x}\]
        积分两次得 $S(x) = (1-x)\ln(1-x) + x, \; x \in (-1, 1)$．\\
        对于端点 $x=1$：由于级数在此处收敛且由 Abel 第二定理知和函数左连续，有
        \[S(1) = \sum_{n=1}^\infty \frac{1}{n(n+1)} = \sum_{n=1}^\infty \left(\frac{1}{n}-\frac{1}{n+1}\right) = 1\]
        而 $\lim\limits_{x \to 1^-} [(1-x)\ln(1-x) + x] = 1$．\\
        综上所述，和函数为
        \[ \boxed{S(x) = \begin{cases} (1-x)\ln(1-x) + x, & -1 \le x < 1 \\ 1, & x = 1 \end{cases}} \]
        \item 易知级数收敛半径 $R=1$．当 $x=\pm 1$ 时级数发散，故成立区间为 $(-1, 1)$．\\
        和函数为
        \[ \boxed{S(x)=\frac{x}{2}\sum_{n=1}^{\infty}\frac{(x^2)^n}{n}=-\frac{x}{2}\ln(1-x^2), \quad x \in (-1, 1)} \]
        \item 易知级数收敛半径 $R=+\infty$，故收敛域为 $(-\infty, +\infty)$．\\
        注意到
        \[S(x) = \sum_{n=0}^\infty \frac{x^{3n}}{(3n)!}, \quad S'(x) = \sum_{n=1}^\infty \frac{x^{3n-1}}{(3n-1)!}, \qquad S''(x) = \sum_{n=1}^\infty \frac{x^{3n-2}}{(3n-2)!} \]
        因此
        \[S(x) + S'(x) + S''(x) = \sum_{n=0}^\infty \frac{x^n}{n!} = e^x\]
        由常微分方程的解法知
        \[ \boxed{S(x) = \dfrac{1}{3} e^x + \dfrac{2}{3} e^{-\frac{x}{2}} \cos\left(\frac{\sqrt{3}x}{2}\right), \quad x \in (-\infty, +\infty) } \]
    \end{enumerate}
\end{solution}

\begin{example}
    将$\ln (4+x)$展开为$x$的幂级数，并指出成立区间．
\end{example}
\begin{solution}
    将函数化为已知形式：$\ln (4+x) = \ln\left[4(1+\frac{x}{4})\right] = \ln 4 + \ln(1+\frac{x}{4})$．\\
    已知 $\ln(1+u) = \displaystyle\sum_{n=1}^\infty (-1)^{n-1} \frac{u^n}{n}, u \in (-1, 1]$．\\
    代入 $u = \frac{x}{4}$，得 \[ \boxed{\ln(4+x) = \ln 4 + \sum_{n=1}^\infty (-1)^{n-1} \frac{x^n}{n \cdot 4^n}} \]
    成立区间为 $-1 < \frac{x}{4} \le 1$，即 $\boxed{x \in (-4, 4]}$．
\end{solution}

\begin{example}
    将函数$f(x)=\dfrac{1}{x^2-3x-4}$展开成$x-1$的幂级数，并指出成立区间．
\end{example}
\begin{solution}
    分解部分分式：$f(x) = \dfrac{1}{(x-4)(x+1)} = \dfrac{1}{5} \left( \dfrac{1}{x-4} - \dfrac{1}{x+1} \right)$．\\
    令 $t = x-1$，则原式为 $\frac{1}{5} \left( \dfrac{1}{t-3} - \dfrac{1}{t+2} \right) = \dfrac{1}{5} \left( -\dfrac{1}{3}\cdot\dfrac{1}{1-\frac{t}{3}} - \dfrac{1}{2}\cdot\dfrac{1}{1+\frac{t}{2}} \right)$．\\
    利用 $\dfrac{1}{1-u} = \displaystyle\sum_{n=0}^\infty u^n$，得：
    \[ f(t) = \frac{1}{5} \left( -\frac{1}{3}\sum_{n=0}^\infty \left(\frac{t}{3}\right)^n - \frac{1}{2}\sum_{n=0}^\infty \left(-\frac{t}{2}\right)^n \right) \]
    回代 $x$：\[ \boxed{f(x) = -\frac{1}{15} \sum_{n=0}^\infty \frac{(x-1)^n}{3^n} - \frac{1}{10} \sum_{n=0}^\infty \frac{(-1)^n (x-1)^n}{2^n}} \]
    收敛域取并集，故 $x-1 \in (-2, 2)$，即 $\boxed{x \in (-1, 3)}$．
\end{solution}

\section{傅里叶级数}
\begin{theorem}[Dirichlet 收敛定理]
    设 $f(x)$ 是周期为 $2\pi$ 的周期函数，在一个周期内连续或只有有限个第一类间断点，并且至多只有有限个极值点，则 $f(x)$ 的傅里叶级数收敛于：
    \[ \frac{a_0}{2} + \sum_{n=1}^\infty (a_n \cos nx + b_n \sin nx) = \frac{f(x+0) + f(x-0)}{2} \]
    其中：
    \[ a_n = \frac{1}{\pi} \int_{-\pi}^{\pi} f(x) \cos nx \d x, \quad b_n = \frac{1}{\pi} \int_{-\pi}^{\pi} f(x) \sin nx \d x \]
\end{theorem}

\begin{corollary}[周期为 $T=2L$ 的傅里叶级数]
    若 $f(x)$ 是周期为 $T=2L$ 的周期函数，且满足 Dirichlet 收敛条件，则其傅里叶级数展开式为：
    \[ f(x) \sim \frac{a_0}{2} + \sum_{n=1}^\infty \left( a_n \cos \frac{n\pi x}{L} + b_n \sin \frac{n\pi x}{L} \right) \]
    其中系数计算公式为：
    \[ a_n = \frac{1}{L} \int_{-L}^{L} f(x) \cos \frac{n\pi x}{L} \d x, \quad (n=0, 1, 2, \dots) \]
    \[ b_n = \frac{1}{L} \int_{-L}^{L} f(x) \sin \frac{n\pi x}{L} \d x, \quad (n=1, 2, \dots) \]
\end{corollary}

\begin{example}
    设周期函数$f(x)$的周期为$2\pi,\,f(x)$在$[-\pi,\pi)$上表达式为$f(x)=e^{2x}$，试将$f(x)$展开为傅里叶级数．
\end{example}
\begin{solution}
    直接套用公式计算系数：
    \[a_0 = \frac{1}{\pi} \int_{-\pi}^\pi e^{2x} \d x = \frac{e^{2\pi} - e^{-2\pi}}{2\pi}\]
    对于 $n \ge 1$，利用分部积分法可得：
    \[a_n = \frac{1}{\pi} \int_{-\pi}^\pi e^{2x} \cos nx \d x = \frac{2(-1)^n(e^{2\pi} - e^{-2\pi})}{\pi(4+n^2)}\]
    \[b_n = \frac{1}{\pi} \int_{-\pi}^\pi e^{2x} \sin nx \d x = \frac{-n(-1)^n(e^{2\pi} - e^{-2\pi})}{\pi(4+n^2)}\]
    故 $f(x)$ 的傅里叶级数为：
    \[ \boxed{ \frac{e^{2\pi} - e^{-2\pi}}{4\pi} + \frac{e^{2\pi} - e^{-2\pi}}{\pi} \sum_{n=1}^\infty \frac{(-1)^n}{4+n^2} (2\cos nx - n\sin nx) } \]
\end{solution}

\begin{example}
    将$f(x)=|\sin x|\;(x\in(-\pi,\pi])$展开成傅里叶级数
\end{example}
\begin{solution}
    $f(x)=|\sin x|$ 是偶函数，故 $b_n = 0$，只需计算 $a_n$．
    \[ a_0 = \frac{2}{\pi} \int_{0}^{\pi} \sin x \d x = \frac{2}{\pi}[-\cos x]_{0}^{\pi} = \frac{4}{\pi} \]
    对于 $n \ge 1$，利用积化和差公式：
    \[ a_n = \frac{2}{\pi} \int_{0}^{\pi} \sin x \cos nx \d x = \frac{1}{\pi} \int_{0}^{\pi} [\sin((n+1)x) - \sin((n-1)x)] \d x \]
    \begin{itemize}
        \item 当 $n=1$ 时：$\sin((n-1)x) = 0$，故
        \[ a_1 = \frac{1}{\pi} \int_0^\pi \sin 2x \d x = \frac{1}{\pi} \left[-\frac{\cos 2x}{2}\right]_0^\pi = 0 \]
        \item 当 $n \neq 1$ 时：
        \begin{align*}
            a_n &= \frac{1}{\pi} \left[ -\frac{\cos((n+1)x)}{n+1} + \frac{\cos((n-1)x)}{n-1} \right]_{0}^{\pi} \\
            &= \frac{1}{\pi} \left[ \left( -\frac{(-1)^{n+1}}{n+1} + \frac{(-1)^{n-1}}{n-1} \right) - \left( -\frac{1}{n+1} + \frac{1}{n-1} \right) \right]
        \end{align*}
        当 $n$ 为奇数 $(n=2k+1,\;k\ge 1)$ 时，$a_{2k+1}=0$；当 $n$ 为偶数 $(n=2k)$ 时：
        \[ a_{2k} = \frac{1}{\pi} \left( -\frac{4}{4k^2-1} \right) = -\frac{4}{\pi(4k^2-1)} \]
    \end{itemize}
    因此 $f(x) = |\sin x|$ 的傅里叶级数展开式为：
    \[ \boxed{|\sin x| = \frac{2}{\pi} - \frac{4}{\pi} \sum_{n=1}^{\infty} \frac{\cos 2nx}{4n^2-1},\quad x\in(-\infty,+\infty)} \]
    由于 $|\sin x|$ 在 $\R$ 上连续且逐段光滑，该展开式在整个实数轴上成立．
\end{solution}

\begin{myRemark}
    在上题中，我们将$2\pi$作为展开的周期，但$|\sin x|$的最小正周期为$\pi$，那么，为什么我们不以$\pi$作为展开的周期呢？\\
    事实上，无论哪一种周期展开，最终的傅里叶级数都是等价的，不过三角函数的频率不同而已。
    但是，在\textbf{应试}中，往往直接采用给定的定义域作为展开的周期，多数答案也是如此。
\end{myRemark}

\begin{example}
    将$f(x)=2+x\;(x\in[0,1])$展开成以$2$为周期的余弦级数，并求$\displaystyle\sum_{n=1}^{\infty}\frac{1}{n^2}$．
\end{example}
\begin{solution}
    余弦级数需对 $f(x)$ 作偶延拓，周期 $T=2L=2$ 即 $L=1$，$b_n=0$．
    \[ a_0 = \frac{2}{L} \int_0^L f(x) \d x = 2 \int_0^1 (2+x) \d x = 2\left[ 2x + \frac{x^2}{2} \right]_0^1 = 5 \]
    对于 $n \ge 1$：
    \begin{align*}
        a_n &= \frac{2}{L} \int_0^L f(x) \cos\frac{n\pi x}{L} \d x = 2 \int_0^1 (2+x) \cos(n\pi x) \d x \\
        &= 2 \left[ (2+x)\frac{\sin(n\pi x)}{n\pi} \Big|_0^1 - \int_0^1 \frac{\sin(n\pi x)}{n\pi} \d x \right] \\
        &= -\frac{2}{n\pi} \int_0^1 \sin(n\pi x) \d x = -\frac{2}{n\pi}\cdot\frac{1-(-1)^n}{n\pi} \\
        &= \frac{2[(-1)^n-1]}{n^2\pi^2} = \begin{cases}
            0, & n \text{ 为偶数},\\
            -\dfrac{4}{n^2\pi^2}, & n \text{ 为奇数}
        \end{cases}
    \end{align*}
    故余弦级数展开式为：
    \[ \boxed{2+x = \frac{5}{2} - \frac{4}{\pi^2} \sum_{k=0}^{\infty} \frac{\cos[(2k+1)\pi x]}{(2k+1)^2},\quad x\in[0,1]} \]
    为求 $\displaystyle\sum_{n=1}^{\infty}\frac{1}{n^2}$，取 $x=0$ 代入上式，得 $f(0)=2$：
    \[ 2 = \frac{5}{2} - \frac{4}{\pi^2} \sum_{k=0}^{\infty} \frac{1}{(2k+1)^2} \]
    \[ \sum_{k=0}^{\infty} \frac{1}{(2k+1)^2} = \frac{\pi^2}{8} \]
    又 $\displaystyle\sum_{n=1}^{\infty}\frac{1}{n^2} = \sum_{k=0}^{\infty}\frac{1}{(2k+1)^2} + \sum_{k=1}^{\infty}\frac{1}{(2k)^2} = \frac{\pi^2}{8} + \frac{1}{4}\sum_{n=1}^{\infty}\frac{1}{n^2}$，解得：
    \[ \boxed{\sum_{n=1}^{\infty}\frac{1}{n^2} = \frac{\pi^2}{6}} \]
\end{solution}

\section{习题}
\begin{exercise}
    $a_n < b_n$，且 $\displaystyle\sum_{n=1}^{\infty}a_n,\sum_{n=1}^{\infty}b_n$ 均收敛，则“$\displaystyle\sum_{n=1}^{\infty}a_n$绝对收敛”是“$\displaystyle\sum_{n=1}^{\infty}b_n$绝对收敛”的（~~~~~~）
    \begin{tasks}[style = itemize, label=\Alph*.](2)
        \task 充分不必要条件
        \task 必要不充分条件
        \task 充分必要条件
        \task 既非充分又非必要条件
    \end{tasks}
\end{exercise}
\begin{solution}
    记 $c_n = b_n - a_n$，由条件 $a_n < b_n$ 知 $c_n > 0$．又 $\sum a_n$ 与 $\sum b_n$ 均收敛，故
    \[ \sum_{n=1}^{\infty} c_n = \sum_{n=1}^{\infty} (b_n - a_n) = \sum_{n=1}^{\infty} b_n - \sum_{n=1}^{\infty} a_n \]
    收敛，即 $\sum c_n$ 是收敛的正项级数．

    \textbf{充分性：} 若 $\sum a_n$ 绝对收敛，则 $\sum |a_n|$ 收敛．由三角不等式得
    \[ |b_n| = |a_n + c_n| \le |a_n| + c_n \]
    从而
    \[ \sum_{n=1}^{\infty} |b_n| \le \sum_{n=1}^{\infty} |a_n| + \sum_{n=1}^{\infty} c_n < +\infty \]
    故 $\sum b_n$ 也绝对收敛．

    \textbf{必要性：} 若 $\sum b_n$ 绝对收敛，则 $\sum |b_n|$ 收敛．由三角不等式得
    \[ |a_n| = |b_n - c_n| \le |b_n| + c_n \]
    从而
    \[ \sum_{n=1}^{\infty} |a_n| \le \sum_{n=1}^{\infty} |b_n| + \sum_{n=1}^{\infty} c_n < +\infty \]
    故 $\sum a_n$ 也绝对收敛．

    综上所述，$\sum a_n$ 绝对收敛 $\iff$ $\sum b_n$ 绝对收敛，即为\textbf{充分必要条件}，选 \boxed{\text{C}}．
\end{solution}

\begin{exercise}
    求下列级数的收敛域：
    \begin{tasks}[label=(\arabic*)](2)
        \task $\displaystyle\sum_{n=1}^{\infty}\left(\frac{a^n}{n}+\frac{b^n}{n^2}\right)x^n\;(a>0,b>0)$
        \task $\displaystyle\sum_{n=0}^{\infty}\frac{x^{n^2}}{2^n}$
    \end{tasks}
\end{exercise}
\begin{solution}
    \begin{enumerate}[label=(\arabic*)]
        \item 设 $u_n = \dfrac{a^n}{n}+\dfrac{b^n}{n^2}$，由根值判别法求收敛半径：
        \[ \sqrt[n]{u_n} = \sqrt[n]{\frac{a^n}{n} + \frac{b^n}{n^2}} \]
        由于 $\sqrt[n]{\dfrac{a^n}{n}} = \dfrac{a}{\sqrt[n]{n}} \to a$，$\sqrt[n]{\dfrac{b^n}{n^2}} = \dfrac{b}{\sqrt[n]{n^2}} \to b$，故
        \[ \lim_{n \to \infty} \sqrt[n]{u_n} = \max(a,b) \]
        从而收敛半径 $R = \dfrac{1}{\max(a,b)} = \dfrac{1}{M},\;M = \max(a,b)$．
        \begin{itemize}
            \item $x = \dfrac{1}{M}$ 时：
            \[ \sum_{n=1}^\infty \left(\frac{a^n}{n}+\frac{b^n}{n^2}\right)\frac{1}{M^n} = \sum_{n=1}^\infty \left[\frac{1}{n}\!\left(\frac{a}{M}\right)^{\!n} + \frac{1}{n^2}\!\left(\frac{b}{M}\right)^{\!n}\right] \]
            \begin{itemize}
                \item 若 $a \ge b$，则 $M = a$，此时 $\frac{a}{M} = 1$，$\frac{b}{M} = \frac{b}{a} < 1$，故第一项为 $\sum \frac{1}{n}$ 发散，整体发散．
                \item 若 $b > a$，则 $M = b$，此时 $\frac{a}{M} = \frac{a}{b} < 1$，$\frac{b}{M} = 1$，故第一项 $\sum \frac{(a/b)^n}{n}$ 收敛（比值法），第二项 $\sum \frac{1}{n^2}$ 收敛，整体收敛．
            \end{itemize}
            \item $x = -\dfrac{1}{M}$ 时：
            \[ \sum_{n=1}^\infty \left(\frac{a^n}{n}+\frac{b^n}{n^2}\right)\frac{(-1)^n}{M^n} = \sum_{n=1}^\infty \frac{(-1)^n}{n}\!\left(\frac{a}{M}\right)^{\!n} + \sum_{n=1}^\infty \frac{(-1)^n}{n^2}\!\left(\frac{b}{M}\right)^{\!n} \]
            \begin{itemize}
                \item 若 $a \ge b$，则 $\frac{a}{M} = 1$，第一项为交错调和级数（条件收敛），第二项绝对收敛，整体收敛．
                \item 若 $b > a$，则 $\frac{a}{M} = \frac{a}{b} < 1$，第一项绝对收敛，第二项亦绝对收敛，整体收敛．
            \end{itemize}
        \end{itemize}
        综上所述，当 $a \ge b$ 时收敛域为 $\boxed{\left[-\dfrac{1}{a},\dfrac{1}{a}\right)}$；当 $a < b$ 时收敛域为 $\boxed{\left[-\dfrac{1}{b},\dfrac{1}{b}\right]}$．

        \item 记 $u_n(x) = \dfrac{x^{n^2}}{2^n}$，由根值判别法：
        \[ \sqrt[n]{|u_n(x)|} = \sqrt[n]{\frac{|x|^{n^2}}{2^n}} = \frac{|x|^n}{2} \]
        当 $|x| < 1$ 时，$\dfrac{|x|^n}{2} \to 0 < 1$，级数收敛；
        当 $|x| = 1$ 时，$\dfrac{|x|^n}{2} = \dfrac{1}{2} < 1$，级数收敛；
        当 $|x| > 1$ 时，$\dfrac{|x|^n}{2} \to +\infty > 1$，级数发散．
        故收敛域为 $|x| \le 1$，即 $\boxed{[-1,\,1]}$．
    \end{enumerate}
\end{solution}
\begin{exercise}
    设 $u_n = \displaystyle\int_{n\pi}^{(n+1)\pi} \frac{\sin x}{\sqrt{1+x^p}} \d x$，讨论级数 $\displaystyle\sum_{n=1}^{\infty} u_n$ 的敛散性．
\end{exercise}
\begin{solution}
    令 $x = t + n\pi$，则 $\d x = \d t$，且 $\sin x = \sin(t+n\pi) = (-1)^n \sin t$，因此
    \[ u_n = (-1)^n \int_0^{\pi} \frac{\sin t}{\sqrt{1+(t+n\pi)^p}} \d t \]
    记 $\displaystyle v_n = \int_0^{\pi} \frac{\sin t}{\sqrt{1+(t+n\pi)^p}} \d t > 0$，则 $\displaystyle\sum_{n=1}^{\infty} u_n = \sum_{n=1}^{\infty} (-1)^n v_n$ 为交错级数．

    \begin{enumerate}[label=\Roman*.]
        \item\textbf{$p \le 0$ 时}
        \begin{itemize}
            \item 若 $p = 0$，则 $v_n = \dfrac{1}{\sqrt{2}} \displaystyle\int_0^{\pi} \sin t \d t = \sqrt{2}$，$\lim\limits_{n\to\infty} |u_n| = \sqrt{2} \neq 0$，级数发散．
            \item 若 $p < 0$，令 $q = -p > 0$，则 $(t+n\pi)^p = (t+n\pi)^{-q} \to 0\;(n\to\infty)$，故
            \[ \lim_{n\to\infty} v_n = \int_0^{\pi} \sin t \d t = 2 \neq 0 \]
            通项不趋于 $0$，级数发散．
        \end{itemize}

        \item\textbf{$p > 0$ 时}
        \begin{itemize}
            \item 单调性：$p > 0$ 时分母 $\sqrt{1+(t+n\pi)^p}$ 随 $n$ 严格递增，故 $v_{n+1} < v_n$．
            \item 极限为零：$0 < v_n \le \displaystyle\int_0^{\pi} \frac{\sin t}{\sqrt{1+(n\pi)^p}} \d t = \frac{2}{\sqrt{1+(n\pi)^p}} \to 0$，由夹逼定理得 $\lim\limits_{n\to\infty} v_n = 0$．
        \end{itemize}
        故 $\sum u_n$ 收敛．

        \item\textbf{$p > 0$ 时绝对收敛性的判定．} 考察正项级数 $\sum v_n$．
        \[ \frac{2}{\sqrt{1+((n+1)\pi)^p}} \le v_n \le \frac{2}{\sqrt{1+(n\pi)^p}} \]
        取比较基准 $\dfrac{1}{n^{\frac{p}{2}}}$，由于
        \[ \lim_{n\to\infty} \frac{v_n}{\frac{1}{n^{\frac{p}{2}}}} = \frac{2}{\pi^{\frac{p}{2}}} \in (0, +\infty) \]
        故 $\sum v_n$ 与 $\sum \dfrac{1}{n^{\frac{p}{2}}}$ 同敛散．由 $p$-级数结论：
        \begin{itemize}
            \item 当 $p > 2$ 时 $\sum \dfrac{1}{n^{\frac{p}{2}}}$ 收敛，从而 $\sum |u_n|$ 收敛，原级数\textbf{绝对收敛}；
            \item 当 $0 < p \le 2$ 时 $\sum \dfrac{1}{n^{\frac{p}{2}}}$ 发散，从而 $\sum |u_n|$ 发散，原级数\textbf{条件收敛}．
        \end{itemize}
    \end{enumerate}

    综上所述，
    \[ \boxed{\begin{cases}
        p > 2, & \text{绝对收敛} \\[4pt]
        0 < p \le 2, & \text{条件收敛} \\[4pt]
        p \le 0, & \text{发散}
    \end{cases}} \]
\end{solution}
\ifx\allfiles\undefined
\end{document}
\fi